% =====================================================================
%  2bat-ud1-14.tex  —  SESSIÓ 14: Repàs global i síntesi (estacions)
%  (sense preàmbul: s'inclou des de main.tex)
%  Criteris: C1, C2, C3, C4. Repàs/síntesi de tota la unitat.
% =====================================================================
\horatitol{Sessió 14 — Repàs global i síntesi}%
{Un recorregut per tota la unitat en quatre estacions: construir i llegir, sumar i escalar, multiplicar, i interpretar.}

\subsection{Mapa de la unitat}

Tota la Unitat 1 es pot resumir en quatre passos encadenats, que són justament les
quatre estacions d'aquest repàs:

\begin{idea}
\textbf{Els quatre blocs de la unitat}
\begin{itemize}[leftmargin=1.4em, itemsep=2pt]
  \item \textbf{A. Organitzar dades (C1):} posar la informació en una matriu,
        llegir-ne la dimensió i localitzar elements.
  \item \textbf{B. Combinar i escalar (C2):} sumar o restar matrices de la mateixa
        dimensió i multiplicar per un escalar (increments i retallades globals).
  \item \textbf{C. Encreuar amb el producte (C3):} multiplicar dues matrices fila
        per columna per obtenir totals (hores, costos).
  \item \textbf{D. Decidir i interpretar (C4):} comprovar la compatibilitat,
        anticipar la dimensió del resultat i llegir-ne el significat.
\end{itemize}
\textit{Aquesta eina obre la porta a la programació lineal de la propera unitat.}
\end{idea}

% ---------------------------------------------------------------------
\subsection{Estació A — Construir i llegir matrices}

\begin{exercicis}[Estació A]
\begin{exlist}
  \item Un centre anota les places ocupades de tres serveis (reforç, lleure,
        menjador) en dos mesos. Al gener: $20$, $15$ i $30$; al febrer: $24$, $18$
        i $28$. Organitza les dades en una matriu (files: serveis; columnes: mesos),
        indica'n la dimensió i digues quin servei té més ocupació al febrer.
  \item A la matriu $A=\begin{pmatrix} 7 & 2 & 5 \\ 1 & 8 & 4 \end{pmatrix}$,
        digues la dimensió, l'element de la fila 2, columna 1, i si és quadrada.
\end{exlist}
\end{exercicis}

\subsection{Estació B — Sumar i escalar}

\begin{exercicis}[Estació B]
\begin{exlist}
  \item Donades $A=\begin{pmatrix} 6 & 3 \\ 2 & 5 \end{pmatrix}$ i
        $B=\begin{pmatrix} 1 & 4 \\ 3 & 0 \end{pmatrix}$, calcula $A+B$ i $A-B$.
  \item Un pressupost és $P=\begin{pmatrix} 2000 & 800 \\ 1500 & 600 \end{pmatrix}$.
        S'aplica una retallada lineal del $10\%$. Per quin escalar cal multiplicar?
        Troba el nou pressupost.
\end{exlist}
\end{exercicis}

\subsection{Estació C — El producte de matrices}

\begin{exercicis}[Estació C]
\begin{exlist}
  \item Calcula $\begin{pmatrix} 2 & 1 \\ 3 & 0 \end{pmatrix}\cdot\begin{pmatrix} 4 & 2 \\ 1 & 5 \end{pmatrix}$.
  \item Un servei atén dos usuaris amb tres tipus de sessió:
        $S=\begin{pmatrix} 2 & 1 & 3 \\ 1 & 4 & 0 \end{pmatrix}$, i les durades són
        $H=\begin{pmatrix} 2 \\ 1 \\ 1 \end{pmatrix}$. Calcula $S\cdot H$ (les hores
        de cada usuari).
\end{exlist}
\end{exercicis}

\subsection{Estació D — Compatibilitat i interpretació}

\begin{exercicis}[Estació D]
\begin{exlist}
  \item Digues si es poden fer, i amb quina dimensió: (a) $(2\times 3)\cdot(3\times 2)$;
        \;(b) $(2\times 2)\cdot(3\times 2)$.
  \item Una matriu d'usuaris per serveis és $4\times 2$ i una de serveis per cost
        és $2\times 1$. Es pot fer el producte? Quina dimensió tindria i què
        representaria el resultat?
\end{exlist}
\end{exercicis}

% ---------------------------------------------------------------------
\fullsolucions{Sessió 14}
\begin{solucions}
\textbf{Estació A}
\begin{sollist}
  \item $\begin{pmatrix} 20 & 24 \\ 15 & 18 \\ 30 & 28 \end{pmatrix}$,
        \;dimensió $3\times 2$. Al febrer (columna 2) el menjador té més ocupació
        ($28$).
  \item Dimensió $2\times 3$; \;l'element de la fila 2, columna 1 és $1$; \;no és
        quadrada.
\end{sollist}

\textbf{Estació B}
\begin{sollist}
  \item $A+B=\begin{pmatrix} 7 & 7 \\ 5 & 5 \end{pmatrix}$,
        \;$A-B=\begin{pmatrix} 5 & -1 \\ -1 & 5 \end{pmatrix}$.
  \item Cal multiplicar per l'escalar $0{,}9$;
        \;nou pressupost $=\begin{pmatrix} 1800 & 720 \\ 1350 & 540 \end{pmatrix}$.
\end{sollist}

\textbf{Estació C}
\begin{sollist}
  \item $\begin{pmatrix} 9 & 9 \\ 12 & 6 \end{pmatrix}$.
  \item $S\cdot H=\begin{pmatrix} 8 \\ 6 \end{pmatrix}$
        (l'usuari 1 genera $8$ hores; l'usuari 2, $6$).
\end{sollist}

\textbf{Estació D}
\begin{sollist}
  \item (a) es pot fer, dimensió $2\times 2$; \;(b) no es pot fer ($2\neq 3$).
  \item Sí: $4\times 2$ per $2\times 1$ dona $4\times 1$. Representaria, per a cada
        usuari (files), el cost total (única columna).
\end{sollist}
\end{solucions}
